Author\+: Julien Lesgourgues


\begin{DoxyItemize}
\item {\bfseries{For what the code can actually compute}}\+: all possible input parameters, all coded cosmological models, all functionalities, all observables, etc.\+: read the file {\ttfamily explanatory.\+ini} in the main {\ttfamily CLASS} directory\+: it is THE reference file where we keep track of all possible input and the definition of all input parameters. For that reason we recommend to leave it always unchanged and to work with copies of it, or with short input files written from scratch.
\item {\bfseries{For the structure, style, and concrete aspects of the code}}\+: this documentation, especially the {\ttfamily CLASS overview} chapter (the extensive automatically-\/generated part of this documentation is more for advanced users); plus the slides of our {\ttfamily CLASS} lectures, for instance those from Tokyo 2014 available at

{\ttfamily \href{http://lesgourg.github.io/class-tour-Tokyo.html}{\texttt{ http\+://lesgourg.\+github.\+io/class-\/tour-\/\+Tokyo.\+html}}}

or the more recent and concise summary from the Narbonne 2016 lecture available at

{\ttfamily \href{http://lesgourg.github.io/class-tour/Narbonne.pdf}{\texttt{ http\+://lesgourg.\+github.\+io/class-\/tour/\+Narbonne.\+pdf}}}

An updated overview of available {\ttfamily CLASS} lecture slides is always available at

{\ttfamily \href{http://lesgourg.github.io/courses.html}{\texttt{ http\+://lesgourg.\+github.\+io/courses.\+html}}}

in the section {\ttfamily Courses on numerical tools}.
\item {\bfseries{For the python wrapper of {\ttfamily CLASS}}}\+: at the moment, the best are the last slides (pages 75-\/96) of the Narbonne 2016 lectures

{\ttfamily \href{http://lesgourg.github.io/class-tour/Narbonne.pdf}{\texttt{ http\+://lesgourg.\+github.\+io/class-\/tour/\+Narbonne.\+pdf}}}

Later we will expand the wrapper documentation with a dedicated chapter here.
\item {\bfseries{For the physics and equations used in the code}}\+: mainly, the following papers\+:
\begin{DoxyItemize}
\item {\itshape Cosmological perturbation theory in the synchronous and conformal Newtonian gauges}

C. P. Ma and E. Bertschinger.

\href{http://arxiv.org/abs/astro-ph/9506072}{\texttt{ http\+://arxiv.\+org/abs/astro-\/ph/9506072}}

10.\+1086/176550

Astrophys. J. {\bfseries{455}}, 7 (1995)
\item {\itshape The Cosmic Linear Anisotropy Solving System (CLASS) II\+: Approximation schemes}

D. Blas, J. Lesgourgues and T. Tram.

\href{http://arxiv.org/abs/1104.2933}{\texttt{ http\+://arxiv.\+org/abs/1104.\+2933}} \mbox{[}astro-\/ph.\+CO\mbox{]}

10.\+1088/1475-\/7516/2011/07/034

JCAP {\bfseries{1107}}, 034 (2011)
\item {\itshape The Cosmic Linear Anisotropy Solving System (CLASS) IV\+: efficient implementation of non-\/cold relics}

J. Lesgourgues and T. Tram.

\href{http://arxiv.org/abs/1104.2935}{\texttt{ http\+://arxiv.\+org/abs/1104.\+2935}} \mbox{[}astro-\/ph.\+CO\mbox{]}

10.\+1088/1475-\/7516/2011/09/032

JCAP {\bfseries{1109}}, 032 (2011)
\item {\itshape Optimal polarisation equations in FLRW universes}

T. Tram and J. Lesgourgues.

\href{http://arxiv.org/abs/1305.3261}{\texttt{ http\+://arxiv.\+org/abs/1305.\+3261}} \mbox{[}astro-\/ph.\+CO\mbox{]}

10.\+1088/1475-\/7516/2013/10/002

JCAP {\bfseries{1310}}, 002 (2013)
\item {\itshape Fast and accurate CMB computations in non-\/flat FLRW universes}

J. Lesgourgues and T. Tram.

\href{http://arxiv.org/abs/1312.2697}{\texttt{ http\+://arxiv.\+org/abs/1312.\+2697}} \mbox{[}astro-\/ph.\+CO\mbox{]}

10.\+1088/1475-\/7516/2014/09/032

JCAP {\bfseries{1409}}, no. 09, 032 (2014)
\item {\itshape The CLASSgal code for Relativistic Cosmological Large Scale Structure}

E. Di Dio, F. Montanari, J. Lesgourgues and R. Durrer.

\href{http://arxiv.org/abs/1307.1459}{\texttt{ http\+://arxiv.\+org/abs/1307.\+1459}} \mbox{[}astro-\/ph.\+CO\mbox{]}

10.\+1088/1475-\/7516/2013/11/044

JCAP {\bfseries{1311}}, 044 (2013)
\end{DoxyItemize}

plus also some latex notes on specific sectors\+:
\begin{DoxyItemize}
\item {\itshape Equations for perturbed recombination}

(can be turned on optionally by the user since v2.\+1.\+0)

L. Voruz.

\href{http://lesgourg.github.io/class_public/perturbed_recombination.pdf}{\texttt{ http\+://lesgourg.\+github.\+io/class\+\_\+public/perturbed\+\_\+recombination.\+pdf}}
\item {\itshape PPF formalism in Newtonian and synchronous gauge}

(used by default for the fluid perturbations since v2.\+6.\+0)

T. Tram.

\href{http://lesgourg.github.io/class_public/PPF_formalism.pdf}{\texttt{ http\+://lesgourg.\+github.\+io/class\+\_\+public/\+PPF\+\_\+formalism.\+pdf}} 
\end{DoxyItemize}
\end{DoxyItemize}